\documentclass{article}
\usepackage[utf8]{inputenc}
\usepackage{amsmath}
\usepackage{amssymb}

\title{Core Mathematical Model}
\date{}

\begin{document}

\maketitle

\section*{Multiresource Flexible Job-Shop Scheduling}

\paragraph{Source and data definitions.}
The selected application is the simultaneous-occupation multi-resource flexible
job-shop problem, PDF pp.2--3 (printed pp.2323--2324), Sections 2 and 3.1,
source equations (1)--(11). Let $J$ be jobs, $\Omega=\bigcup_{u\in J}O_u$
be their operations, and $R$ be actual unary resources. Operation $i$ has
requirement groups $G_i$; group $j$ must receive one resource from
$R_{i,j}\subseteq R$. The same actual resource cannot serve two groups of
the same operation. $p_{i,k}\ge0$ are fixed processing times. $PJ_i$ are
immediate predecessors in a given acyclic precedence graph. All assigned
resources are occupied simultaneously for the common operation duration.
$\mathcal M$ is a valid disjunctive constant for the finite scheduling horizon;
a feasible serial schedule supplies a makespan upper bound when needed.
The independent variables are nonnegative starts $\tau_i$, durations $p_i$,
makespan $C_{\max}\ge0$, assignments $Y$, resource-use indicators $U$ and
pairwise order indicators $X$, with the binary domains below.
The source completion variable is $t_i=\tau_i+p_i$; this changes notation,
defined by its assignment meaning. The source's printed equation (3)
has a double sum inconsistent with its explanation: we use its explicitly
stated per-resource rule, also confirmed by CP constraint (14). The duration
epigraph is retained as in source (4); at an optimum durations may be reduced
to the largest selected resource processing time without worsening makespan.


For practical Gurobi implementation, we use start-time variables $\tau_i$, duration variables $p_i$, and auxiliary binary variables $U_{i,k}$ indicating whether operation $i$ uses resource $k$. Let $R_i = \bigcup_{j \in G_i} R_{i,j}$ and let $\mathcal{P} = \{(i,i') \in \Omega \times \Omega : i < i'\}$. This is equivalent to the paper's MILP but avoids ambiguous time notation.

\begin{equation}
\min C_{\max}
\tag{1}
\end{equation}

\begin{equation}
\sum_{k \in R_{i,j}} Y_{i,j,k} = 1,
\qquad
\forall i \in \Omega,\ \forall j \in G_i
\tag{2}
\end{equation}

\begin{equation}
\sum_{j \in G_i:\ k \in R_{i,j}} Y_{i,j,k} \le 1,
\qquad
\forall i \in \Omega,\ \forall k \in R
\tag{3}
\end{equation}

\begin{equation}
U_{i,k} =
\sum_{j \in G_i:\ k \in R_{i,j}} Y_{i,j,k},
\qquad
\forall i \in \Omega,\ \forall k \in R
\tag{4}
\end{equation}

\begin{equation}
p_i \ge \sum_{k \in R_{i,j}} p_{i,k} Y_{i,j,k},
\qquad
\forall i \in \Omega,\ \forall j \in G_i
\tag{5}
\end{equation}

\begin{equation}
\tau_i \ge \tau_{i'} + p_{i'},
\qquad
\forall i \in \Omega,\ \forall i' \in PJ_i
\tag{6}
\end{equation}

\begin{equation}
\tau_{i'} \ge
\tau_i + p_i
- \mathcal{M}\left(3 - X_{i,i',k} - U_{i,k} - U_{i',k}\right),
\quad
\forall (i,i') \in \mathcal{P},\ \forall k \in R_i \cap R_{i'}
\tag{7}
\end{equation}

\begin{equation}
\tau_i \ge
\tau_{i'} + p_{i'}
- \mathcal{M}\left(2 + X_{i,i',k} - U_{i,k} - U_{i',k}\right),
\quad
\forall (i,i') \in \mathcal{P},\ \forall k \in R_i \cap R_{i'}
\tag{8}
\end{equation}

\begin{equation}
C_{\max} \ge \tau_i + p_i,
\qquad \forall i \in \Omega
\tag{9}
\end{equation}

\begin{equation}
\tau_i \ge 0,\quad p_i \ge 0,
\qquad \forall i \in \Omega
\tag{10}
\end{equation}

\begin{equation}
X_{i,i',k} \in \{0,1\},
\qquad \forall (i,i') \in \mathcal{P},\ \forall k \in R_i \cap R_{i'}
\tag{11}
\end{equation}

\begin{equation}
Y_{i,j,k} \in \{0,1\},
\qquad \forall i \in \Omega,\ \forall j \in G_i,\ \forall k \in R_{i,j}
\tag{12}
\end{equation}

\begin{equation}
U_{i,k} \in \{0,1\},
\qquad \forall i \in \Omega,\ \forall k \in R
\tag{13}
\end{equation}

\end{document}
